The linear representation hypothesis holds that large language models (LLMs) encode concepts as linear directions in activation space, yet recent work reveals that some features---notably refusal behavior---exhibit fundamentally nonlinear, multi-dimensional structure. We ask whether this geometric distinction predicts a model's ability to introspect on its own internal states, drawing an analogy to the conscious/unconscious divide in human cognition: linearly encoded features may be ``accessible'' to self-report, while nonlinearly encoded features may not. We test this \emph{Introspection--Linearity Hypothesis} through a two-pronged design. First, we extract residual-stream activations from \qwen and train linear and nonlinear probes to quantify the linearity of three feature types: sentiment, binary refusal, and harmful-content sub-categories. Second, we administer introspection tests to \gptfour, measuring self-report accuracy on tasks of varying linearity. We find that perfectly linear features (sentiment, binary refusal) yield 100\% self-report accuracy, while tasks involving nonlinear structure show degraded performance: sub-category classification drops to 97.5\% and borderline refusal prediction to 90.0\%. Confidence calibration reveals a significant gap between clear and ambiguous cases ($p = 0.021$). These results provide partial support for the hypothesis that representation geometry constrains introspective access in LLMs, with implications for alignment monitoring and interpretability.
